\documentclass[a4paper]{article}
\usepackage{amsfonts}
\usepackage{amsmath}
\usepackage{amssymb}
\usepackage{graphicx}     
\newcommand{\ddx}{\dfrac{d}{dx}}
\newcommand{\ddg}{\dfrac{d}{dg}}
\newcommand{\dint}{\displaystyle\int}
\newcommand{\bgtan}{\operatorname{Bgtan}}

\begin{document}
  
\noindent \large Auteur: Luca Letroye \\
\noindent \large Studierichting: Informatica\\
\noindent \large Oefeningenreeks 4A nr. 45.2

\medskip
\normalsize

$\dint\dfrac{18x^2}{x^4-4x+4}dx = \dint\dfrac{18x^2}{(x-1)^2(x^2+2x+3)}dx$\\

$ \dfrac{18x^2}{(x-1)^2(x^2+2x+3)} = \dfrac{A}{(x-1)^2} + \dfrac{B}{(x-1)} + \dfrac{Cx+D}{x^2+2x+3}$\\

$A = \left.\dfrac{18x^2}{x^2+2x+3}\right|_{x=1} = 3$\\

$\dfrac{18x^2}{(x-1)^2(x^2+2x+3)} - \dfrac{3}{(x-1)^2} = \dfrac{3(5x+3)}{(x-1)(x^2+2x+3)}$\\

$\Rightarrow B = \left.\dfrac{3(5x+3)}{x^2+2x+3}\right|_{x=1} = 4$\\

$18x^2 = 3(x^2+2x+3) + B(x-1)(x^2+2x+3) (Cx+D)(x-1)^2$\\

$\Rightarrow 0x^3+4x^3+Cx^3 \Rightarrow C = -4$\\

$x = 0 \Rightarrow D = 3$\\

$\Rightarrow 3\dint\dfrac{dx}{(x-1)^2}+4\dint\dfrac{dx}{x-1}+\dint\dfrac{-4x+3}{x^2+2x+3}dx$\\

$ = 3\dint t^{-2}dt +4\ln|x-1| - \dint\dfrac{2(2x+2)}{x^2+2x+3}+\dint\dfrac{7}{x^2+2x+3}dx $\\

$\dint\dfrac{7}{x^2+2x+3}dx = 7\dint\dfrac{1}{(x+1)^2+(\sqrt{2})^2} = \dfrac{7}{\sqrt{2}}\bgtan\left(\dfrac{x+1}{\sqrt{2}}\right)+C$\\

$\Rightarrow \dfrac{-3}{x-1}+4\ln|x-1| -2\ln|x^2+2x+3| +\dfrac{7}{\sqrt{2}}\bgtan\left(\dfrac{x+1}{\sqrt{2}}\right)+C$\\



\begin{thebibliography}{999}
%\bibitem{a} Referentie
\end{thebibliography}
\end{document} 
